\documentclass[conference,a4paper]{IEEEtran}
\usepackage[T1]{fontenc}
\usepackage[utf8]{inputenc}
\usepackage[spanish,es-tabla]{babel}
\usepackage{lmodern}
\usepackage{microtype}
\usepackage{csquotes}
\usepackage{booktabs,tabularx,array,xurl}
\usepackage{amsmath}
\usepackage[hidelinks]{hyperref}
\usepackage[backend=biber,style=ieee,sorting=none,maxbibnames=3]{biblatex}
\addbibresource{referencias.bib}
\DefineBibliographyStrings{spanish}{bibliography={Referencias}}
\newcolumntype{Y}{>{\raggedright\arraybackslash}X}
\newcolumntype{P}[1]{>{\raggedright\arraybackslash}p{#1}}
\setlength{\parindent}{0.7em}
\setlength{\parskip}{1pt}
\setlength{\columnsep}{0.24in}
\renewcommand{\arraystretch}{1.18}
\urlstyle{same}
\title{Proceso de requisitos aplicado a una plataforma de ventas\\con recomendaciones basadas en IA para establecimientos el\'ectricos}
\author{\IEEEauthorblockN{Cede\~no Avila Winston Damian\\
Paca Pilataza Luis Alberto\\
Rivera Suarez Jonathan Javier}
\IEEEauthorblockA{Universidad T\'ecnica Estatal de Quevedo (UTEQ)\\
Ingenier\'ia de Software, cuarto nivel ``B'' --- ISR-401\\
Pr\'actica Experimental Unidad I --- GA-SUM-01}}
\begin{document}
\maketitle
\begin{abstract}
Esta versi\'on es un avance progresivo de la pr\'actica experimental de Ingenier\'ia de Requisitos (ISR-401). Se aplica el enfoque de ISO/IEC/IEEE 29148:2018, Pohl y SWEBOK a un PFC de ventas con recomendaciones de IA para comercios de materiales el\'ectricos. Las necesidades, responsabilidades y conexiones descritas son propuestas para revisar con el negocio: no se presentan entrevistas ni pruebas de sistema como si ya hubieran ocurrido.
\end{abstract}
\section*{0. Selecci\'on del PFC y contexto del problema}
\addcontentsline{toc}{section}{0. Selecci\'on del PFC y contexto del problema}
\subsection*{Descripci\'on del sistema (m\'aximo 150 palabras)}
El proyecto propone una plataforma de ventas para establecimientos que comercializan materiales el\'ectricos. Su objetivo es reunir en un mismo entorno el registro de productos, la consulta de disponibilidad, la preparaci\'on de ventas y el historial de operaciones. Como apoyo, plantea un m\'odulo de recomendaci\'on basado en inteligencia artificial que sugiera art\'iculos relacionados con la solicitud del comprador, sin sustituir la verificaci\'on t\'ecnica ni la decisi\'on del vendedor. Los usuarios previstos son personal de ventas, administradores y responsables de inventario. El escenario de referencia del PFC es un establecimiento del sector el\'ectrico de Quevedo (ELECTRORIOS); sus procedimientos concretos, permisos de acceso y fuentes de informaci\'on a\'un deben confirmarse con los responsables. Esta propuesta se analiza como un sistema por construir, no como un producto desplegado ni evaluado en condiciones reales.

\subsection*{Problema y finalidad}
En una venta de material el\'ectrico, el trabajador debe identificar la necesidad del cliente y comprobar si el producto existe, tiene un precio vigente y resulta pertinente. Si esa informaci\'on queda dispersa, se incrementa el riesgo de consultas repetidas o sugerencias poco apropiadas. La oportunidad de software es reunir el proceso comercial y aportar recomendaciones \emph{explicables}, siempre sujetas a revisi\'on humana. Este enfoque se formula como una necesidad preliminar que debe contrastarse con quienes usan el negocio; una necesidad no equivale a un requisito aprobado \cite{iso29148,pohl2025}.

\subsection*{Alcance inicial y l\'imites}
\begin{itemize}
\item \textbf{Dentro del alcance propuesto:} consulta de cat\'alogo, stock disponible, registro de ventas, perfiles de acceso y sugerencias de productos con justificaci\'on visible.
\item \textbf{Fuera del alcance, por ahora:} comprobaci\'on autom\'atica de compatibilidad el\'ectrica sin un especialista, decisiones aut\'onomas de compra, y conexiones con servicios externos cuyo contrato no haya sido confirmado.
\item \textbf{No verificado:} inventario real del establecimiento, frecuencia de compras, herramientas actuales, conectividad y funcionamiento de cualquier motor de IA. Todo esto requiere levantamiento autorizado.
\end{itemize}

\subsection*{Tipo de sistema y visi\'on preliminar}
Se clasifica de manera preliminar como un sistema de informaci\'on transaccional con funciones de apoyo a decisiones de venta. El componente de IA ampl\'ia las consultas, pero no reemplaza la responsabilidad del personal ni convierte las recomendaciones en validaciones t\'ecnicas de seguridad. La separaci\'on entre procesos del producto y reglas organizacionales ayudar\'a a evitar requisitos imposibles de comprobar \cite{pohl2025,iso29148}.

\textbf{Visi\'on (adaptaci\'on al formato Moore de la gu\'ia):} Para el personal de establecimientos el\'ectricos que necesita consultar art\'iculos y completar ventas con informaci\'on coherente, la plataforma es un sistema de ventas con apoyo de recomendaci\'on que centraliza datos de productos y propone opciones relacionadas. A diferencia de llevar registros aislados, busca ofrecer trazabilidad del proceso y sugerencias que el vendedor pueda revisar.

\subsection*{Necesidades preliminares por validar}
\begin{table}[!ht]
\centering\footnotesize
\caption{Necesidades iniciales que deben contrastarse con el establecimiento}
\label{tab:necesidades}
\begin{tabularx}{\columnwidth}{@{}P{0.65cm}Y P{1.25cm}@{}}
\toprule ID & Necesidad propuesta & Por validar\\\midrule
N-01 & Consultar precio y stock sin duplicar registros. & Observaci\'on\\
N-02 & Registrar la venta y recuperar su referencia. & Entrevista\\
N-03 & Contar con sugerencias de art\'iculos pertinentes y revisables. & Taller\\
N-04 & Restringir modificaciones de precios y existencias por rol. & Aprobaci\'on\\
\bottomrule
\end{tabularx}
\end{table}
\noindent\textit{Estado:} estas necesidades son hip\'otesis de trabajo. Las t\'ecnicas citadas son acciones previstas; no se han realizado entrevistas ni pruebas comerciales.

\subsection*{Preguntas para la primera elicitaci\'on}
\begin{enumerate}
\item \textquestiondown{}C\'omo se identifica un art\'iculo y se consulta su disponibilidad antes de venderlo?
\item \textquestiondown{}Qui\'en est\'a autorizado a cambiar precios, existencias y datos de clientes?
\item \textquestiondown{}Qu\'e informaci\'on debe conservarse para revisar una venta?
\item \textquestiondown{}Cu\'ando ser\'ia \`util recomendar un producto alternativo y qui\'en decide si es apto?
\item \textquestiondown{}Cu\'ales son los problemas m\'as frecuentes en el flujo comercial actual?
\end{enumerate}

\section{Fundamento te\'orico aplicado al PFC}
\subsection{2.1. Ingenier\'ia de Requisitos y necesidades}
La Ingenier\'ia de Requisitos (IR) permite transformar necesidades de personas y organizaciones en acuerdos que gu\'ien la construcci\'on de un sistema. No comienza redactando funciones por intuici\'on: identifica el contexto, distingue qu\'e problema es relevante y documenta requisitos que pueden revisarse. La norma ISO/IEC/IEEE 29148:2018 describe los procesos e informaci\'on necesarios para desarrollar requisitos de sistemas y software a lo largo del ciclo de vida \cite{iso29148}. El temario Foundation Level del IREB, versi\'on 3.1 indicada por la gu\'ia docente, complementa los conceptos de elicitaci\'on y especificaci\'on \cite{ireb2022}. Pohl presenta la IR como una pr\'actica que relaciona obtenci\'on, negociaci\'on, documentaci\'on, validaci\'on y gesti\'on en un marco coherente \cite{pohl2025}.

Para el PFC, pedir ``una aplicaci\'on que venda m\'as'' expresa una aspiraci\'on general, pero no indica qu\'e debe hacer la plataforma. Conviene primero examinar c\'omo el personal atiende una solicitud, consulta un producto y confirma su disponibilidad. A partir de esa informaci\'on pueden identificarse necesidades distintas: evitar registros duplicados, conocer existencias y revisar qu\'e recomendaciones son oportunas. Cada necesidad debe tener una fuente y quedar abierta a correcci\'on antes de convertirse en especificaci\'on.

\subsection{2.2. Niveles y tipos de requisitos}
Los requisitos de negocio representan resultados buscados por una organizaci\'on; los de interesados describen necesidades de grupos afectados; los de usuario permiten reconocer tareas; y los de sistema y software detallan las condiciones que el producto debe cumplir. La distinci\'on impide confundir una meta comercial con una decisi\'on tecnol\'ogica. Las gu\'ias de IR recomiendan mantener relaciones trazables entre niveles y no dar por aprobados supuestos sin responsable \cite{iso29148,swebok2024}.

\textbf{Aplicaci\'on:} una meta de negocio podr\'ia ser disminuir la repetici\'on de consultas de disponibilidad. Como necesidad del vendedor se propone ver en una sola pantalla la descripci\'on del art\'iculo, su precio y su stock. Un requisito funcional preliminar ser\'ia mostrar esos datos cuando el vendedor seleccione un c\'odigo de producto autorizado. Se trata de una propuesta acad\'emica: la meta, las pantallas y sus reglas de acceso requieren contraste con la empresa. No conviene definir desde el inicio una base de datos o una API concreta sin conocer las restricciones del PFC.

\subsection{2.3. Funciones, interfaces, restricciones y calidad}
Un requisito funcional indica un comportamiento comprobable. Por su parte, los requisitos de interfaz especifican intercambios con personas u otros sistemas, las restricciones acotan opciones de soluci\'on y los requisitos de calidad precisan las condiciones bajo las cuales el producto debe funcionar \cite{iso29148}. Glinz advierte que la etiqueta ``no funcional'' no tiene una interpretaci\'on universal y requiere una definici\'on cuidadosa en el proyecto \cite{glinz2007}. Por eso, en este informe se usa la etiqueta ``requisito de calidad'' cuando pueda asociarse con una caracter\'istica del modelo ISO/IEC 25010:2023 \cite{iso25010}.

Por ejemplo, ``el sistema ser\'a r\'apido'' no permite aceptar o rechazar el resultado. Un objetivo provisional m\'as preciso es que el 95\,\% de las consultas al cat\'alogo terminen en tres segundos o menos bajo una carga de ensayo acordada. El valor de tres segundos \emph{no es una medici\'on de ELECTRORIOS}; deber\'a negociarse y probarse. De forma semejante, un requisito de seguridad puede exigir que un vendedor no edite precios sin permiso, y uno de interfaz puede definir qu\'e datos se enviar\'ian al motor de recomendaci\'on sin incluir identificadores personales innecesarios. La necesidad de una interfaz real deber\'a corroborarse antes del dise\~no.

\subsection{2.4. Propiedades de buena formulaci\'on}
Un requisito individual debe ser necesario, apropiado, no ambiguo, completo, singular, factible, verificable, correcto y conforme al estilo acordado. Estas propiedades se discuten en la norma ISO/IEC/IEEE 29148:2018, especialmente en sus apartados 5.2.4--5.2.7 \cite{iso29148}. No basta con que una frase suene t\'ecnica: ``mostrar sugerencias adecuadas'' deja sin definir qu\'e es adecuado. Si el vendedor recibe una alternativa, se deber\'ia aclarar qu\'e informaci\'on acompa\~na la sugerencia y qu\'e acci\'on permite aceptarla o descartarla.

Un conjunto tambi\'en debe ser consistente, comprensible y susceptible de validaci\'on. Si un requisito permite vender sin existencias y otro lo impide, el conflicto debe resolverse con el responsable del proceso antes de programar. Las expresiones ``f\'acil'', ``siempre'' y ``lo m\'as pronto posible'' requieren condiciones medibles o deben reemplazarse. La plantilla EARS ayuda a representar un disparador y una respuesta observable sin mezclar varios comportamientos \cite{mavin2009}. Por ejemplo: ``Cuando el vendedor seleccione un producto activo, el sistema deber\'a mostrar el precio vigente registrado para ese producto''. Las decisiones sobre descuentos, impuestos o stock se a\~nadir\'an tras consultar la pol\'itica real del negocio.

\subsection{Relaci\'on con el marco y ciclo de vida}
Pohl propone considerar el contexto como parte central de la IR, y SWEBOK organiza el conocimiento del \`area de requisitos junto con otros campos de ingenier\'ia de software \cite{pohl2025,swebok2024}. La gu\'ia del docente tambi\'en solicita ISO/IEC/IEEE 15288:2023 para procesos de ciclo de vida de sistemas y la edici\'on ISO/IEC/IEEE 12207:2017 para software \cite{iso15288,iso12207}. Esta \`ultima se conserva aqu\'i como \emph{edici\'on solicitada en la gu\'ia}, aunque ISO ha publicado una versi\'on posterior. La selecci\'on definitiva de procesos se har\'a al conocer los acuerdos y recursos disponibles del PFC.

\noindent\textit{Nota de investigaci\'on:} las fichas normativas y los datos editoriales citados se verificaron en los portales de ISO, IEEE, Springer y IEEE Computer Society. Para una cita literal o interpretaci\'on m\'as espec\'ifica de una cl\'ausula, el equipo deber\'a contrastar la norma completa mediante el acceso institucional.

\section*{Alcance del avance}
Adem'as del contexto, se document'o el fundamento te'orico inicial 2.1--2.4, con ejemplos que deben discutirse en el PFC. Contin'uan pendientes los resultados de campo y la evaluaci'on individual de requisitos.
\section*{Referencias}
\printbibliography[heading=none]
\noindent\footnotesize\textbf{Declaraci\'on de uso de IA.} Se utiliz\'o ChatGPT (OpenAI, GPT-6) para preparar estos borradores en LaTeX, ordenar el contenido y apoyar la revisi\'on bibliogr\'afica. Se contrastaron los datos editoriales con fichas oficiales de ISO, IEEE, Springer, PMI e IREB. Los integrantes deben leer, ajustar y validar sus aportes; las observaciones, entrevistas y acuerdos no realizados no se presentan como evidencias.
\end{document}
